\documentclass[11pt,a4paper]{beamer}
\usetheme{Madrid}
\usepackage[utf8]{inputenc}
\usepackage{amsmath}
\usepackage{amsfonts}
\usepackage{amssymb}
\usepackage{graphicx}
\usepackage{xcolor}
\usepackage{tikz}
\usepackage{booktabs}
\usepackage{array}
\usepackage{colortbl}
\usepackage{multicol}

% ============================================================
% EXTREMELY SMALL FONT - MAXIMUM CONTENT PER SLIDE
% ============================================================
\usefonttheme{professionalfonts}

% Keep things compact but stable
\setbeamerfont{normal text}{size=\scriptsize}
\setbeamerfont{itemize/enumerate body}{size=\scriptsize}
\setbeamerfont{itemize/enumerate subbody}{size=\tiny}
\setbeamerfont{block body}{size=\scriptsize}
\setbeamerfont{block title}{size=\scriptsize}
\setbeamerfont{frametitle}{size=\small}
\setbeamerfont{title}{size=\normalsize}
\setbeamerfont{subtitle}{size=\small}
\setbeamerfont{author}{size=\tiny}
\setbeamerfont{date}{size=\tiny}

\renewcommand{\scriptsize}{\fontsize{6pt}{7pt}\selectfont}
\renewcommand{\footnotesize}{\fontsize{6.5pt}{8pt}\selectfont}
\renewcommand{\small}{\fontsize{7.5pt}{9pt}\selectfont}
\renewcommand{\normalsize}{\fontsize{8.5pt}{10.5pt}\selectfont}

\newcommand{\redflag}[1]{\textcolor{red}{\textbf{[!] #1}}}
\newcommand{\bluenote}[1]{\textcolor{blue}{\textbf{[i] #1}}}
\newcommand{\greennote}[1]{\textcolor{green!50!black}{\textbf{[Right] #1}}}
\newcommand{\examfocus}[1]{\textcolor{orange!80!black}{\textbf{[FOCUS] #1}}}
\newcommand{\trap}[1]{\textcolor{purple}{\textbf{[TRAP] #1}}}
\newcommand{\correct}[1]{\textcolor{green!50!black}{\textbf{[CORRECT] #1}}}
\newcommand{\scenario}[1]{\textcolor{blue!70!black}{\textbf{[SCENARIO] #1}}}
\newcommand{\keyterm}[1]{\textcolor{red!80!black}{\textbf{[KEY] #1}}}
\newcommand{\lawcite}[1]{\textcolor{brown!70!black}{\textbf{[LAW] #1}}}

\title[CAMS Domain A]{CAMS Exam Preparation\\Domain A: Risks and Methods of\\Money Laundering and Terrorist Financing (20 percent)}
\author{Advanced AML Training Framework}
\date{}

\setbeamercovered{transparent}
\setbeamertemplate{navigation symbols}{}
\setbeamertemplate{frametitle continuation}{}
\setlength{\parskip}{2pt}
\setlength{\itemsep}{1pt}

\begin{document}
	
	% ============================================================
	% TITLE SLIDE
	% ============================================================
	\begin{frame}
		\titlepage
		\centering
		\tiny Domain A: 20 percent of CAMS Exam | 100+ Slides | ML Stages, TBML, PEPs, NPOs, Human Trafficking, Virtual Assets
	\end{frame}
	
	% ============================================================
	% SECTION 1: THREE STAGES OF MONEY LAUNDERING
	% ============================================================
	
	\begin{frame}{A.1: The Three Stages of Money Laundering}
		\begin{block}{Core Concept}
			Money laundering is the process of making criminal proceeds appear legitimate through three distinct stages.
		\end{block}
		\begin{tabular}{p{0.25\textwidth}|p{0.35\textwidth}|p{0.35\textwidth}}
			\hline
			\rowcolor{lightgray}
			\textbf{Stage} & \textbf{Definition} & \textbf{Examples} \\
			\hline
			\textbf{Placement} & Introducing illicit funds into the financial system & Cash deposits, currency smuggling, purchasing assets with cash \\
			\hline
			\textbf{Layering} & Conducting complex transactions to obscure the audit trail & Wire transfers through multiple accounts, shell companies, trade invoices \\
			\hline
			\textbf{Integration} & Making laundered funds appear legitimate & Investing in real estate, luxury goods, business acquisitions \\
			\hline
		\end{tabular}
		\vspace{0.2cm}
		\redflag{Exam tests ability to identify which stage is occurring in a scenario.}
		\examfocus{Key distinction: Placement = first entry into system. Layering = obscuring the trail. Integration = making funds appear legitimate.}
	\end{frame}
	
	\begin{frame}{A.1.1: Placement Stage - Detailed Analysis}
		\begin{block}{Placement Methods}
			Placement is the most vulnerable stage for criminals because large amounts of cash are difficult to introduce without detection.
		\end{block}
		\textbf{Common Placement Techniques:}
		\begin{itemize}
			\item Cash deposits into bank accounts (often structured below CTR thresholds)
			\item Currency smuggling across borders (cash couriers)
			\item Purchasing monetary instruments (money orders, traveler's checks)
			\item Using casinos (buying chips with cash, minimal gambling, cashing out)
			\item Commingling with legitimate business revenues (cash-intensive businesses)
			\item Purchasing high-value assets (cars, art, jewelry) with cash
			\item Using prepaid cards (load with cash, use or resell)
			\item Hawala and informal value transfer systems
		\end{itemize}
		\textbf{Red Flags for Placement:}
		\begin{itemize}
			\item Cash deposits just below reporting thresholds (structuring)
			\item Multiple deposits from different locations on same day
			\item Cash-intensive business with deposits inconsistent with expected volume
			\item Customer unable to explain source of cash
		\end{itemize}
		\redflag{Structuring (smurfing) is a federal crime independent of the underlying source of funds.}
	\end{frame}
	
	\begin{frame}{A.1.2: Layering Stage - Detailed Analysis}
		\begin{block}{Layering Methods}
			Layering creates distance between the criminal proceeds and their origin through complex transactions.
		\end{block}
		\textbf{Common Layering Techniques:}
		\begin{itemize}
			\item Multiple wire transfers through different banks and jurisdictions
			\item Using shell companies and offshore accounts
			\item Converting funds to different currencies
			\item Purchasing and selling assets (real estate, securities)
			\item Using trade-based money laundering (over/under-invoicing)
			\item Using virtual assets and mixers/tumblers
			\item Using casinos (chip washing)
			\item Using online gaming platforms and in-game currencies
		\end{itemize}
		\textbf{Red Flags for Layering:}
		\begin{itemize}
			\item Rapid movement of funds through multiple accounts
			\item Round-dollar wire transfers with no apparent business purpose
			\item Transfers to/from high-risk jurisdictions or secrecy havens
			\item Use of multiple intermediaries with no logical connection
			\item Transactions involving shell companies with opaque ownership
		\end{itemize}
		\redflag{Layering is the most complex stage and often involves multiple jurisdictions.}
	\end{frame}
	
	\begin{frame}{A.1.3: Integration Stage - Detailed Analysis}
		\begin{block}{Integration Methods}
			Integration returns laundered funds to the criminal as apparently legitimate wealth.
		\end{block}
		\textbf{Common Integration Techniques:}
		\begin{itemize}
			\item Purchasing legitimate businesses (restaurants, retail stores, real estate)
			\item Investing in securities and financial instruments
			\item Using funds as collateral for legitimate loans
			\item Creating false invoices and business records
			\item Reporting laundered funds as business income on tax returns
			\item Purchasing luxury assets (yachts, art, aircraft)
			\item Establishing trusts and foundations
		\end{itemize}
		\textbf{Red Flags for Integration:}
		\begin{itemize}
			\item Purchase of high-value assets through shell companies
			\item Business with no clear source of funding for expansion
			\item Real estate purchases with all-cash through anonymous LLCs
			\item Unexplained wealth inconsistent with known income
			\item Transactions involving previously flagged layering accounts
		\end{itemize}
		\redflag{Integration makes laundered funds appear legitimate - often the hardest stage to detect.}
	\end{frame}
	
	\begin{frame}{A.1.4: Three-Stage Analysis - Exam Scenario}
		\scenario{Multi-Stage Laundering Scheme}
		An organized crime group generates \$4 million in narcotics proceeds. The group's accountant deposits cash at 14 different branches, keeping each deposit between \$7,500 and \$9,800. The funds are consolidated into a holding LLC, which issues fictitious management-fee invoices to the original businesses. The LLC uses the funds to purchase three residential rental properties. Rental income is reported to the IRS as legitimate business income.
		
		\vspace{0.1cm}
		\textbf{What the Exam Tests:} Identifying which transaction corresponds to which stage.
		
		\vspace{0.1cm}
		\textbf{How to Analyze:}
		\begin{enumerate}
			\item Step 1 (cash deposits): PLACEMENT - introducing cash into system
			\item Step 2 (fictitious invoices): LAYERING - creating false paper trail
			\item Step 3 (property purchase): INTEGRATION - investing in legitimate assets
			\item Step 4 (rental income): INTEGRATION - making funds appear legitimate
		\end{enumerate}
		
		\trap{Common Distractor: "The fictitious invoices represent integration."}
		\correct{Correct Answer: Fictitious invoices create a paper trail to obscure origin = LAYERING. Property purchase with laundered funds = INTEGRATION.}
	\end{frame}
	
	% ============================================================
	% SECTION 2: TRADE-BASED MONEY LAUNDERING (TBML)
	% ============================================================
	
	\begin{frame}{A.2: Trade-Based Money Laundering (TBML)}
		\begin{block}{What is TBML?}
			The process of disguising criminal proceeds through trade transactions using over-invoicing, under-invoicing, phantom shipments, or multiple invoicing.
		\end{block}
		\textbf{Primary TBML Techniques:}
		\begin{itemize}
			\item Over-invoicing: Exporter inflates invoice price to transfer value from importer to exporter
			\item Under-invoicing: Exporter lowers invoice price to move value from exporter to importer
			\item Phantom shipping: No goods actually shipped; only paper documentation exists
			\item Multiple invoicing: Same goods invoiced multiple times to different buyers
			\item Short-shipping: Shipping less than invoiced amount
			\item Misrepresentation of goods: False description of goods on customs forms
		\end{itemize}
		\textbf{Red Flags for TBML:}
		\begin{itemize}
			\item Invoice price significantly above or below market value
			\item Multiple amendments to letters of credit without logical explanation
			\item Round-dollar invoices (not typical for genuine trade)
			\item Shipment routing through free trade zones or unusual transshipment points
			\item Inconsistent cargo description, weight, or value
			\item Shell company involvement on either side of transaction
		\end{itemize}
		\redflag{TBML is considered the most difficult ML typology to detect.}
	\end{frame}
	
	\begin{frame}{A.2.1: Over-Invoicing vs. Under-Invoicing - Exam Critical}
		\begin{block}{Over-Invoicing - Value Transfer to Exporter}
			Over-invoicing moves value from the importing country to the exporting country.
		\end{block}
		\textbf{Over-Invoicing Example:}
		Market price = \$180 per unit. Invoice price = \$1,200 per unit. Importer pays \$1,200, but goods are only worth \$180. The extra \$1,020 per unit is transferred to the exporter disguised as legitimate trade payment.
		
		\begin{block}{Under-Invoicing - Value Transfer to Importer}
			Under-invoicing moves value from the exporting country to the importing country.
		\end{block}
		\textbf{Under-Invoicing Example:}
		Market price = \$400 per unit. Invoice price = \$80 per unit. Importer pays \$80 but can sell goods at \$400 market price. The \$320 per unit difference becomes laundered cash for the importer.
		
		\trap{Common Distractor: "Under-invoicing benefits the exporter."}
		\correct{Correct Answer: Over-invoicing transfers value TO exporter; Under-invoicing transfers value TO importer.}
	\end{frame}
	
	\begin{frame}{A.2.2: TBML Exam Scenario - Under-Invoicing}
		\scenario{Grossly Under-Invoiced Electronics Shipment}
		A trade finance bank reviews a letter of credit for export of 5,000 industrial electronics units. Invoice price = \$80 per unit. Global market price = \$400 per unit. Exporter is a shell company in a secrecy haven incorporated 11 days ago. Importer is a licensed manufacturer in a FATF-compliant country. Payment routed through the exporter's account at a shell bank.
		
		\vspace{0.1cm}
		\textbf{What the Exam Tests:} Recognizing under-invoicing as a TBML technique.
		
		\vspace{0.1cm}
		\textbf{How to Analyze:}
		\begin{enumerate}
			\item Market price is \$400, invoice is \$80 = 80 percent under-invoicing
			\item Value is being transferred from exporter to importer
			\item Importer will sell goods at market price, generating laundered cash
			\item The exporter's shell company and secrecy haven jurisdiction are additional red flags
		\end{enumerate}
		
		\trap{Common Distractor: "The exporter is over-invoicing to move value out of the country."}
		\correct{Correct Answer: This is UNDER-INVOICING. The importer pays less but sells at market price, generating laundered legitimate cash.}
	\end{frame}
	
	\begin{frame}{A.2.3: Phantom Shipping - Fictitious Trade}
		\scenario{Phantom Shipment with No Physical Goods}
		A bank reviews a trade finance transaction for "industrial machinery components" from a recently incorporated supplier with no verifiable manufacturing capability. Invoice value is \$2.5 million. Shipping documents include a bill of lading from a carrier that, upon investigation, has no record of the shipment. The same supplier name appears on three other pending transactions from different buyers.
		
		\vspace{0.1cm}
		\textbf{What the Exam Tests:} Recognizing phantom shipping red flags.
		
		\vspace{0.1cm}
		\textbf{How to Analyze:}
		\begin{enumerate}
			\item No actual goods exist - only paper documentation
			\item Supplier has no manufacturing capability
			\item Bill of lading is fraudulent (carrier has no record)
			\item Same supplier appearing across multiple suspicious transactions is critical red flag
			\item Bank should decline financing and file SAR
		\end{enumerate}
		
		\trap{Common Distractor: "All documentary requirements appear valid, so the bank should process the transaction."}
		\correct{Correct Answer: Phantom shipping uses fabricated documentation to move funds. The bank should decline, file SAR, and consider terminating the buyer relationship.}
	\end{frame}
	
	\begin{frame}{A.2.4: Free Trade Zone Vulnerabilities}
		\begin{block}{FTZ Money Laundering Risks}
			Free Trade Zones have relaxed regulatory oversight, minimal customs inspections, and intense volume of re-exported goods.
		\end{block}
		\textbf{FTZ Red Flags:}
		\begin{itemize}
			\item Cargo passing through FTZ with no legitimate economic reason
			\item Altered bills of lading showing different origin after FTZ processing
			\item Uniform invoice values regardless of actual cargo type
			\item Phantom shipments (containers declared as goods but found empty)
			\item Mislabeled cargo (luxury electronics declared as low-value scrap)
			\item Nominee directors with no legitimate business role
		\end{itemize}
		\textbf{Exam Scenario Example:}
		A shipping company operating exclusively in a Free Trade Zone processed 340 containers declared as "industrial parts" with uniform invoice value of \$30,000. Physical inspection revealed one container empty, one containing household furniture, and one containing luxury electronics worth \$400,000. Bills of lading altered to show generic descriptions.
		
		\trap{Common Distractor: "FTZs are supervised by customs authorities, so no additional AML scrutiny is needed."}
		\correct{Correct Answer: FTZs create TBML vulnerabilities including phantom shipping and cargo mislabeling. Banks must apply enhanced scrutiny to FTZ transactions.}
	\end{frame}
	
	% ============================================================
	% SECTION 3: SHELL COMPANIES AND BENEFICIAL OWNERSHIP
	% ============================================================
	
	\begin{frame}{A.3: Shell Companies and Money Laundering}
		\begin{block}{What Makes Shell Companies Vulnerable?}
			Shell companies have no physical presence, no employees, and no legitimate business operations - making them ideal for layering.
		\end{block}
		\textbf{Shell Company Red Flags:}
		\begin{itemize}
			\item Use of bearer shares or nominee shareholders to obscure ownership
			\item Multiple layers of corporate entities across secrecy havens
			\item Registered agent address is a mail drop or corporate service provider
			\item No verifiable business purpose or operations
			\item Recently incorporated with no operating history
			\item Transactions inconsistent with stated business purpose
			\item Refusal to provide beneficial ownership information
		\end{itemize}
		\textbf{Bearer Shares Vulnerability:}
		Bearer shares transfer ownership through physical delivery of the stock certificate - no public registry. The true owner is anonymous. FATF requires countries to either prohibit bearer shares or require immobilization with a custodian.
		
		\redflag{Bearer shares + secrecy haven jurisdiction = highest risk.}
	\end{frame}
	
	\begin{frame}{A.3.1: Beneficial Ownership - Aggregating Direct and Indirect Stakes}
		\begin{block}{UBO Calculation Example}
			Under FATF standards, beneficial ownership is calculated by aggregating direct and indirect ownership.
		\end{block}
		\textbf{Example Structure:}
		\begin{itemize}
			\item Individual K owns 18 percent of TargetCo directly
			\item Individual K owns 90 percent of Zenith Holdings LLC
			\item Zenith Holdings LLC owns 15 percent of TargetCo
			\item UBO threshold = 25 percent combined
		\end{itemize}
		\textbf{Calculation:}
		18 percent direct + (90 percent x 15 percent indirect) = 18 percent + 13.5 percent = 31.5 percent.
		
		\vspace{0.1cm}
		\textbf{Result:} Individual K exceeds the 25 percent threshold and MUST be identified and verified as a UBO.
		
		\trap{Common Distractor: "Only direct ownership counts; indirect stakes are irrelevant."}
		\correct{Correct Answer: Both direct and indirect ownership must be aggregated. The calculation multiplies through each ownership tier.}
	\end{frame}
	
	\begin{frame}{A.3.2: Multi-Tier Corporate Structures - Exam Scenario}
		\scenario{Three Layers of Shell Companies}
		An investigation reveals a complex structure: Company Alpha is owned by Company Beta (Panama), which is owned by bearer shares held by a corporate service provider in the Bahamas. The custodian refuses to disclose the bearer share holder's identity. The beneficial owner cannot be identified from any public registry.
		
		\vspace{0.1cm}
		\textbf{What the Exam Tests:} Bearer share risks and CDD obligations.
		
		\vspace{0.1cm}
		\textbf{How to Analyze:}
		\begin{enumerate}
			\item Bearer shares create anonymity for the true owner
			\item Immobilization with a custodian is insufficient if custodian won't disclose identity
			\item The bank cannot complete adequate CDD without UBO identification
			\item The account should be declined
		\end{enumerate}
		
		\trap{Common Distractor: "Immobilization of bearer shares with a licensed custodian satisfies all BO requirements."}
		\correct{Correct Answer: The bank still needs to identify the beneficial owner. Immobilization alone does not eliminate the need for UBO identification.}
	\end{frame}
	
	% ============================================================
	% SECTION 4: BLACK MARKET PESO EXCHANGE (BMPE)
	% ============================================================
	
	\begin{frame}{A.4: Black Market Peso Exchange (BMPE)}
		\begin{block}{What is BMPE?}
			A currency exchange system that uses drug dollars in the U.S. to pay legitimate invoices for foreign merchants, allowing Colombian peso brokers to profit from the exchange rate spread.
		\end{block}
		\textbf{BMPE Mechanism:}
		\begin{enumerate}
			\item Drug cartels have U.S. dollars from narcotics sales
			\item Colombian peso broker buys the dollars at a discount
			\item Broker uses the dollars to pay U.S. invoices of Colombian importers
			\item Colombian importers repay the broker in pesos
			\item Cartel receives pesos in Colombia (cleaned through the trade system)
		\end{enumerate}
		\textbf{Red Flags for BMPE:}
		\begin{itemize}
			\item U.S. importer receives structured cash deposits from unrelated third parties
			\item No logical relationship between depositors and importer
			\item Corresponding imports into Colombia at below-market prices
			\item No direct wire transfers between U.S. and Colombia
			\item Cash deposits used to pay commercial invoices for foreign merchants
		\end{itemize}
		\redflag{BMPE allows drug cartels to convert U.S. dollars to Colombian pesos without cross-border currency movement.}
	\end{frame}
	
	\begin{frame}{A.4.1: BMPE Exam Scenario}
		\scenario{U.S. Importer Receives Structured Cash Deposits}
		A U.S. textile importer's accounts receive \$8.7 million in structured cash deposits from 43 different individuals with no apparent business relationship. Concurrently, a Colombian textile wholesaler receives shipments at 35 percent below wholesale price, paid in pesos by a local broker. No funds are wired directly between Colombia and the United States.
		
		\vspace{0.1cm}
		\textbf{What the Exam Tests:} Recognizing BMPE as the criminal mechanism.
		
		\vspace{0.1cm}
		\textbf{How to Analyze:}
		\begin{enumerate}
			\item Cash deposits from unrelated third parties = narco-dollar collection
			\item Broker uses dollars to pay U.S. importer's invoices
			\item Colombian importer repays broker in pesos
			\item Result: dollars stay in U.S., pesos delivered to cartel in Colombia
			\item No cross-border currency movement = evades detection
		\end{enumerate}
		
		\trap{Common Distractor: "This is a standard hawala arrangement."}
		\correct{Correct Answer: This is BMPE - narco-dollars pay U.S. commercial invoices; Colombian importers repay in pesos. The broker profits from the exchange rate spread.}
	\end{frame}
	
	% ============================================================
	% SECTION 5: POLITICALLY EXPOSED PERSONS (PEPs)
	% ============================================================
	
	\begin{frame}{A.5: Politically Exposed Persons (PEPs)}
		\begin{block}{Definition of PEP}
			Individuals entrusted with prominent public functions, their family members, and close associates.
		\end{block}
		\textbf{Who Qualifies as a PEP?}
		\begin{itemize}
			\item Heads of state and government
			\item Senior politicians and cabinet ministers
			\item Senior judicial officials
			\item Senior military officials
			\item Senior executives of state-owned enterprises
			\item Members of supreme audit institutions
			\item Members of the boards of central banks
			\item Ambassadors and senior diplomats
		\end{itemize}
		\textbf{PEP Family Members and Close Associates:}
		\begin{itemize}
			\item Spouse, partner, children, parents, siblings
			\item Business partners or close advisors
			\item Beneficiaries of trusts or foundations established by the PEP
		\end{itemize}
		\redflag{PEP status requires ENHANCED DUE DILIGENCE and senior management approval.}
	\end{frame}
	
	\begin{frame}{A.5.1: Foreign vs. Domestic PEPs}
		\begin{block}{Key Distinction}
			Foreign PEPs generally require more intensive due diligence than domestic PEPs.
		\end{block}
		\begin{tabular}{p{0.45\textwidth}|p{0.45\textwidth}}
			\hline
			\rowcolor{lightgray}
			\textbf{Foreign PEP} & \textbf{Domestic PEP} \\
			\hline
			Holds prominent position in another country & Holds prominent position in same country \\
			Highest risk - requires full EDD & Moderate risk - EDD required but may be less intensive \\
			Source of wealth and funds verification required & Source of wealth/funds may be more accessible \\
			Senior management approval required & Approval may be delegated \\
			Enhanced ongoing monitoring required & Monitoring commensurate with risk \\
			\hline
		\end{tabular}
		\vspace{0.2cm}
		\textbf{International Organization PEPs:}
		Senior officials of international organizations (UN, IMF, World Bank, EU) are also considered PEPs.
		
		\trap{Common Distractor: "Domestic PEPs are exempt from enhanced due diligence."}
		\correct{Correct Answer: FATF requires enhanced due diligence for domestic PEPs, though intensity may be less than for foreign PEPs.}
	\end{frame}
	
	\begin{frame}{A.5.2: PEP Family Members - Exam Scenario}
		\scenario{Adult Daughter of Foreign Minister Opens Account}
		Ms. Y, adult daughter of the Minister of Energy of an oil-rich nation with high corruption indices, wishes to deposit \$18 million from "consulting fees" paid by energy logistics firms to her BVI company. She holds no government position.
		
		\vspace{0.1cm}
		\textbf{What the Exam Tests:} PEP family member classification and EDD requirements.
		
		\vspace{0.1cm}
		\textbf{How to Analyze:}
		\begin{enumerate}
			\item Ms. Y is a PEP FAMILY MEMBER (not a PEP herself)
			\item PEP family members require Enhanced Due Diligence
			\item The ambiguous "consulting fees" from energy firms related to her father's ministry is a critical red flag
			\item Source of wealth and funds must be verified through corroborative documentation
		\end{enumerate}
		
		\trap{Common Distractor: "Ms. Y holds no government position, so standard CDD is sufficient."}
		\correct{Correct Answer: PEP family members require EDD. The ambiguous consulting fees from her father's ministry sector require intensified scrutiny.}
	\end{frame}
	
	\begin{frame}{A.5.3: Former PEP Grace Period}
		\begin{block}{Post-Retirement PEP Treatment}
			PEP status should be maintained for a reasonable period after retirement, typically at least 12 months but potentially longer based on risk.
		\end{block}
		\textbf{Factors for Ongoing PEP Status:}
		\begin{itemize}
			\item Country's corruption risk level
			\item Former PEP's continued access to state influence
			\item Ongoing business connections to the government
			\item Time elapsed since leaving office
			\item Account activity patterns
		\end{itemize}
		\textbf{Example:}
		A former senior minister retired 5 years ago with 5 years of clean activity. Lowering the risk rating may be justified, but the decision must be documented and risk-based. The bank must consider the country's corruption risk and any ongoing government connections.
		
		\trap{Common Distractor: "A former PEP can never be reclassified to lower risk."}
		\correct{Correct Answer: Former PEPs may be reclassified after reasonable time (typically 12+ months) based on documented risk assessment, including country risk and ongoing influence.}
	\end{frame}
	
	% ============================================================
	% SECTION 6: HUMAN TRAFFICKING FINANCIAL INDICATORS
	% ============================================================
	
	\begin{frame}{A.6: Human Trafficking vs. Human Smuggling}
		\begin{block}{Critical Distinction for Exam}
			Human trafficking and human smuggling are different crimes with different financial indicators.
		\end{block}
		\begin{tabular}{p{0.45\textwidth}|p{0.45\textwidth}}
			\hline
			\rowcolor{lightgray}
			\textbf{Human Trafficking} & \textbf{Human Smuggling} \\
			\hline
			Exploitation (forced labor, sexual exploitation) & Illegal border crossing only \\
			Ongoing, coercive relationship & One-time transaction \\
			Victims controlled by trafficker & Smuggled person is free after crossing \\
			Funds typically move from victim to controller & Funds typically move from migrant to smuggler \\
			Multiple small, regular payments & One-time large payment \\
			Controller pays victim's living expenses & No ongoing relationship \\
			\hline
		\end{tabular}
		\vspace{0.2cm}
		\redflag{Human trafficking = ongoing exploitation. Human smuggling = one-time border crossing.}
	\end{frame}
	
	\begin{frame}{A.6.1: Human Trafficking Financial Red Flags - Forced Labor}
		\scenario{Construction Foreman Account Pattern}
		A construction foreman's account receives 150 peer-to-peer payments ($200-$600 each) from 35 different individuals. Every Friday, 90 percent of the balance is withdrawn in cash. The account pays utility bills for three high-density residential properties. No employment contracts or payroll records exist.
		
		\vspace{0.1cm}
		\textbf{What the Exam Tests:} Recognizing forced labor trafficking indicators.
		
		\vspace{0.1cm}
		\textbf{How to Analyze:}
		\begin{enumerate}
			\item Multiple victims' wages deposited to single controller account
			\item Controller pays utility bills for multiple residential properties (housing victims)
			\item Cash withdrawals suggest physical collection of funds
			\item No employment contracts or payroll records = no legitimate business
		\end{enumerate}
		
		\trap{Common Distractor: "This is a legitimate payroll arrangement for a small construction crew."}
		\correct{Correct Answer: The pattern is highly indicative of human trafficking for forced labor. The analyst should consider filing a SAR.}
	\end{frame}
	
	\begin{frame}{A.6.2: Human Trafficking - Commercial Sexual Exploitation}
		\scenario{Massage Parlor Network Accounts}
		A bank identifies 12 accounts linked to the same group. Each account receives multiple small electronic payments ($50-$150) from dozens of individuals daily. Funds transfer to a central LLC that holds title to six commercial properties zoned for "personal services businesses" (massage parlors). The LLC pays mortgages, utilities, and property taxes. Individuals have no verifiable income other than deposits.
		
		\vspace{0.1cm}
		\textbf{What the Exam Tests:} Recognizing sex trafficking financial indicators.
		
		\vspace{0.1cm}
		\textbf{How to Analyze:}
		\begin{enumerate}
			\item Multiple small payments from many individuals = customer payments for commercial sex
			\item Funds layered through accounts to real estate LLC
			\item Real estate purchase with trafficking proceeds = integration stage
			\item No verifiable income other than deposits = red flag
		\end{enumerate}
		
		\trap{Common Distractor: "Personal services businesses are low-risk customer types."}
		\correct{Correct Answer: The pattern is consistent with human trafficking for commercial sexual exploitation. The bank should file a SAR on the accounts and LLC.}
	\end{frame}
	
	% ============================================================
	% SECTION 7: NON-PROFIT ORGANIZATIONS AND TF
	% ============================================================
	
	\begin{frame}{A.7: NPO Vulnerabilities to Terrorist Financing}
		\begin{block}{Why NPOs Are Vulnerable}
			NPOs enjoy public trust, access high volumes of cash, and routinely operate in or near conflict zones where financial monitoring is weak.
		\end{block}
		\textbf{NPO Terrorist Financing Red Flags:}
		\begin{itemize}
			\item Large donation from a foreign charity operating in a conflict zone
			\item Rapid cash withdrawal (80 percent+ within 48 hours) far from headquarters
			\item Cash withdrawals inconsistent with stated charitable mission
			\item No verifiable charitable activities or documentation
			\item Transfers to other NPOs in conflict zones where terrorist groups operate
			\item Leadership cannot provide documentation of charitable expenditures
			\item NPO with no physical presence, employees, or verifiable activities (shell NPO)
		\end{itemize}
		\textbf{FATF Recommendation 8 Requirements:}
		\begin{itemize}
			\item Targeted risk-based supervision of NPOs vulnerable to TF abuse
			\item Not all NPOs - only those identified as vulnerable
			\item Mechanisms to ensure NPOs are not misused by terrorist organizations
			\item Outreach and awareness programs for NPOs
		\end{itemize}
		\redflag{NPOs operating in conflict zones where terrorist groups control aid distribution require EDD.}
	\end{frame}
	
	\begin{frame}{A.7.1: NPO Terrorist Financing - Exam Scenario}
		\scenario{Conflict Zone Charity with Rapid Cash Withdrawal}
		An NPO registered in a low-risk jurisdiction receives a \$2 million donation from a foreign charity operating in Conflict Zone Z where terrorist groups operate. Within 48 hours, the NPO withdraws \$1.6 million in cash from a branch 300km from its headquarters, in a region with known terrorist activity. The remaining funds are wired to other NPOs in conflict zones. Leadership cannot provide documentation of school construction matching the cash withdrawals.
		
		\vspace{0.1cm}
		\textbf{What the Exam Tests:} Recognizing NPO terrorist financing red flags.
		
		\vspace{0.1cm}
		\textbf{How to Analyze:}
		\begin{enumerate}
			\item Rapid cash withdrawal (80 percent within 48 hours) is a significant red flag
			\item Withdrawal location far from headquarters suggests diversion
			\item Donation from conflict zone with terrorist activity requires enhanced scrutiny
			\item No documentation of charitable activities = funds likely diverted
		\end{enumerate}
		
		\trap{Common Distractor: "NPOs often operate in remote areas where cash is the only practical payment method."}
		\correct{Correct Answer: The rapid cash withdrawal velocity mismatch with charitable mission is the most indicative red flag. Bank should file SAR and consider terminating the account.}
	\end{frame}
	
	% ============================================================
	% SECTION 8: VIRTUAL ASSETS AND VASPs
	% ============================================================
	
	\begin{frame}{A.8: Virtual Asset Money Laundering Risks}
		\begin{block}{Key Virtual Asset Vulnerabilities}
			Virtual assets offer speed, cross-border transferability, and varying degrees of anonymity.
		\end{block}
		\textbf{Virtual Asset Red Flags:}
		\begin{itemize}
			\item Use of anonymity-enhanced cryptocurrencies (Monero, Zcash)
			\item Use of crypto mixers and tumblers to break transaction trail
			\item Rapid conversion between different virtual assets
			\item Use of cross-chain bridges to move between blockchains
			\item Transactions involving DeFi protocols with no KYC
			\item Customer refuses to explain source of virtual assets
			\item Transfer from known illicit wallet (ransomware, darknet markets)
		\end{itemize}
		\textbf{FATF Travel Rule for Virtual Assets:}
		\begin{itemize}
			\item Originator and beneficiary information must be collected and shared
			\item Required for transfers above threshold (typically 1,000 USD/EUR)
			\item Both sending and receiving VASPs have obligations
		\end{itemize}
		\redflag{Privacy coins + mixers + cross-chain bridges = high-risk layering pattern.}
	\end{frame}
	
	\begin{frame}{A.8.1: Privacy Coins and Mixers - Exam Scenario}
		\scenario{Customer Deposits Monero, Uses Mixer, Converts to Bitcoin, Withdraws Fiat}
		A customer deposits 500 Monero (XMR), sends it through a non-regulated mixer, converts to Bitcoin, and requests fiat withdrawal of \$450,000 to a bank in a weak AML jurisdiction. Customer refuses to explain source of XMR, stating "it's private personal savings."
		
		\vspace{0.1cm}
		\textbf{What the Exam Tests:} Recognizing privacy coin and mixer layering patterns.
		
		\vspace{0.1cm}
		\textbf{How to Analyze:}
		\begin{enumerate}
			\item Monero uses ring signatures and stealth addresses - obscures transaction trail
			\item Mixer pools funds and redistributes - breaks audit trail
			\item Rapid conversion to Bitcoin then fiat = layering
			\item Refusal to explain source after mixing is significant red flag
		\end{enumerate}
		
		\trap{Common Distractor: "The use of Monero alone is standard practice and requires no additional scrutiny."}
		\correct{Correct Answer: The sequence of privacy coin deposit, mixing, conversion, and rapid fiat withdrawal is a classic layering pattern requiring EDD and potential SAR filing.}
	\end{frame}
	
	\begin{frame}{A.8.2: NFT Wash Trading - Exam Scenario}
		\scenario{Circular NFT Trading to Inflate Value}
		17 wallets controlled by Individual P buy and sell the same NFTs back and forth 89 times over 60 days, with prices escalating from \$2,000 to \$340,000. All wallets are funded from the same exchange account. Individual P then sells one NFT to an unconnected third party for \$310,000 and converts to fiat.
		
		\vspace{0.1cm}
		\textbf{What the Exam Tests:} Recognizing NFT wash trading as a money laundering technique.
		
		\vspace{0.1cm}
		\textbf{How to Analyze:}
		\begin{enumerate}
			\item Circular trading between self-controlled wallets = wash trading
			\item Artificially inflates price to create fraudulent market record
			\item Selling to unwitting third party at inflated price = layering/integration
			\item Exchange that funded wallets should file SAR
		\end{enumerate}
		
		\trap{Common Distractor: "NFTs are considered art, not financial instruments, so AML regulations don't apply."}
		\correct{Correct Answer: NFT wash trading is a recognized money laundering typology. The exchange should file a SAR based on the circular trading pattern.}
	\end{frame}
	
	% ============================================================
	% SECTION 9: CORRESPONDENT BANKING RISKS
	% ============================================================
	
	\begin{frame}{A.9: Correspondent Banking Money Laundering Risks}
		\begin{block}{Key Correspondent Banking Vulnerabilities}
			Correspondent banking involves one bank (correspondent) providing services to another bank (respondent) - creating potential AML gaps.
		\end{block}
		\textbf{Correspondent Banking Red Flags:}
		\begin{itemize}
			\item Respondent bank services unlicensed MSBs or NBFIs
			\item Nested transactions from undisclosed downstream clients
			\item Large round-dollar wire transfers with vague descriptions
			\item Transactions routed through multiple intermediaries
			\item Respondent bank in jurisdiction with weak AML enforcement
			\item Respondent bank's customer base includes high-risk segments
			\item Payable-through accounts with no CDD on underlying accountholders
		\end{itemize}
		\textbf{Nested Account Risk:}
		When a respondent bank allows third-party banks to route transactions through the correspondent account without disclosure, the correspondent loses visibility into true originators and beneficiaries - a critical due diligence gap.
		
		\redflag{Nested transactions without disclosure = complete loss of visibility.}
	\end{frame}
	
	\begin{frame}{A.9.1: Undisclosed Nested Correspondent - Exam Scenario}
		\scenario{Respondent Bank Routes Transactions for Unlicensed Crypto Exchange}
		A European correspondent bank discovers that its respondent bank has been routing transactions for CryptoX, an unregulated virtual asset exchange, through the correspondent account without disclosure. Transactions are large round-sum wires ($750k-$2 million) with vague descriptions like "consulting fees." Ultimate originators and beneficiaries are not identifiable.
		
		\vspace{0.1cm}
		\textbf{What the Exam Tests:} Correspondent bank obligations for nested accounts.
		
		\vspace{0.1cm}
		\textbf{How to Analyze:}
		\begin{enumerate}
			\item Undisclosed nesting means correspondent lost visibility
			\item This violates core correspondent banking due diligence standards
			\item Correspondent must file SAR and issue RFI to respondent
			\item Relationship may need to be terminated
		\end{enumerate}
		
		\trap{Common Distractor: "The respondent bank holds a valid banking license, so the correspondent has no liability."}
		\correct{Correct Answer: The correspondent bank faces direct regulatory exposure. It must file a SAR, issue an RFI to the respondent, and reassess whether to continue the relationship.}
	\end{frame}
	
	\begin{frame}{A.9.2: Payable-Through Account (PTA) Risks}
		\begin{block}{PTA Definition and Obligations}
			A PTA allows the respondent bank's customers to initiate transactions directly through the correspondent's payment systems.
		\end{block}
		\textbf{Correspondent Bank PTA Obligations (U.S. Regulations):}
		\begin{itemize}
			\item Must apply CDD requirements directly to underlying PTA accountholders
			\item Cannot rely solely on respondent bank's CDD
			\item Must maintain records identifying each PTA accountholder
			\item Must monitor transactions from PTA accountholders
		\end{itemize}
		\textbf{PTA Red Flags:}
		\begin{itemize}
			\item Multiple PTA accountholders sending consistent \$9,500 wires to same shell company
			\item Wires occurring within minutes of each other on same business days
			\item No individual CDD performed on PTA accountholders at onboarding
			\item Aggregate flow exceeds expected volume
		\end{itemize}
		
		\trap{Common Distractor: "The PTA agreement shifts all AML responsibility to the respondent bank."}
		\correct{Correct Answer: Under U.S. regulations, the correspondent bank must apply CDD directly to underlying PTA accountholders.}
	\end{frame}
	
	% ============================================================
	% SECTION 10: CASINO AND GAMBLING AML RISKS
	% ============================================================
	
	\begin{frame}{A.10: Casino Money Laundering Techniques}
		\begin{block}{Casino Vulnerabilities}
			Casinos are used to convert criminal cash into negotiable instruments (checks, chips) with an appearance of legitimate gambling origin.
		\end{block}
		\textbf{Casino Red Flags:}
		\begin{itemize}
			\item Purchasing large volume of chips with cash, minimal gambling, cashing out for check
			\item Multiple individuals purchasing chips for same player
			\item Using third-party checks or wire transfers to purchase chips
			\item Structured cash purchases just below reporting thresholds
			\item Playing low-risk games (e.g., slot machines with high payout rates) to preserve value
			\item "Chip washing" - converting cash to chips and back to check with small loss
		\end{itemize}
		\textbf{Casino AML Obligations (U.S. BSA):}
		\begin{itemize}
			\item Casinos with gross gaming revenue > \$1 million are financial institutions
			\item Must file CTRs for cash transactions exceeding \$10,000 in a gaming day
			\item Must file SARs for suspicious activity exceeding \$5,000
			\item Must have AML compliance programs
		\end{itemize}
		\redflag{"Chip washing" combines placement and integration - cash becomes a documented casino check.}
	\end{frame}
	
	\begin{frame}{A.10.1: Casino Chip Wash - Exam Scenario}
		\scenario{Player Purchases \$300,000 in Chips, Plays 45 Minutes, Cashes Out for Check}
		A player purchases \$300,000 in chips with cash, plays blackjack for 45 minutes with $1,500 in chips, then cashes out $297,000 for a casino check, accepting the \$3,000 loss as a "laundering fee."
		
		\vspace{0.1cm}
		\textbf{What the Exam Tests:} Recognizing chip washing as money laundering.
		
		\vspace{0.1cm}
		\textbf{How to Analyze:}
		\begin{enumerate}
			\item Large cash purchase of chips = placement
			\item Minimal gambling = not legitimate gaming
			\item Cashing out for check = integration (check appears to be gambling winnings)
			\item The \$3,000 loss is the cost of laundering
		\end{enumerate}
		
		\trap{Common Distractor: "The player actually gambled, so the transaction has a legitimate gaming purpose."}
		\correct{Correct Answer: The casino must file a CTR for the \$300,000 cash-in AND a SAR for the chip washing pattern (minimal gambling, cashing out for check).}
	\end{frame}
	
	% ============================================================
	% SECTION 11: HAWALA AND IVTS
	% ============================================================
	
	\begin{frame}{A.11: Informal Value Transfer Systems (IVTS) / Hawala}
		\begin{block}{What is Hawala?}
			An informal value transfer system operating outside formal banking, based on trust and reciprocal obligations between brokers (hawaladars).
		\end{block}
		\textbf{Hawala Characteristics:}
		\begin{itemize}
			\item Transactions settled through trust-based obligations, not formal banking records
			\item Hawaladars maintain informal ledgers (hawala books) to track reciprocal obligations
			\item No physical currency crosses borders
			\item Client gives cash to hawaladar; counterpart delivers equivalent in another country
			\item Settlements occur through future transactions (netting)
		\end{itemize}
		\textbf{Why Hawala is a High AML Risk:}
		\begin{itemize}
			\item No traceable banking record of underlying transaction
			\item Identities of parties may not be recorded
			\item Operates outside formal financial system
			\item Used by legitimate businesses (migrant remittances) AND criminal networks
		\end{itemize}
		\textbf{Regulatory Requirements:}
		\begin{itemize}
			\item FATF Recommendation 14 requires MVTS providers to be licensed or registered
			\item Hawala operators are NOT exempt - must be licensed
			\item Unlicensed operators subject to criminal sanctions
		\end{itemize}
		\redflag{Hawala operators must register as MSBs - no cultural heritage exemption.}
	\end{frame}
	
	\begin{frame}{A.11.1: Hawala Network - Exam Scenario}
		\scenario{Hawaladar Operating from Convenience Store, No MSB Registration}
		A hawaladar in New York accepts \$80,000 in cash from a courier, records a coded reference number, and calls a counterpart in Karachi who pays the equivalent in rupees. No wire transfer occurs; no bank account is used. Over 36 months, the network processed an estimated \$22 million. The hawaladar has no FinCEN MSB registration.
		
		\vspace{0.1cm}
		\textbf{What the Exam Tests:} MSB registration requirements for hawala operators.
		
		\vspace{0.1cm}
		\textbf{How to Analyze:}
		\begin{enumerate}
			\item Hawala is a money or value transfer service (MVTS)
			\item FATF Recommendation 14 requires licensing or registration
			\item Hawaladar is required to register as an MSB with FinCEN
			\item Failure to register is a criminal violation (18 U.S.C. Section 1960)
			\item Processing criminal proceeds adds money laundering charges
		\end{enumerate}
		
		\trap{Common Distractor: "Hawala is exempt from AML regulation due to its cultural heritage status."}
		\correct{Correct Answer: Hawaladar must register as an MSB. Failure to register and processing criminal proceeds exposes him to criminal liability under 18 U.S.C. Section 1960 and money laundering charges.}
	\end{frame}
	
	% ============================================================
	% SECTION 12: MONEY MULES
	% ============================================================
	
	\begin{frame}{A.12: Money Mule Networks}
		\begin{block}{What is a Money Mule?}
			A person who receives and transfers illicit funds through their account, often unknowingly recruited, for a commission.
		\end{block}
		\textbf{Money Mule Red Flags:}
		\begin{itemize}
			\item Account receives multiple incoming P2P transfers from dozens of different individuals
			\item Transfers typically under \$500 each (below reporting thresholds)
			\item Immediate aggregated wire transfers to offshore shell company
			\item Account holder (often student) with no verifiable income
			\item Transfers out shortly after receiving funds
			\item Inconsistent with customer's stated occupation or expected activity
		\end{itemize}
		\textbf{Money Mule Recruitment Typology:}
		\begin{itemize}
			\item "Work from home" job offers
			\item "Transfer funds for a commission" arrangements
			\item Romance scams (victim acts as mule unwittingly)
			\item Student recruitment (easy money for students)
		\end{itemize}
		\redflag{The mule is often a victim of a recruitment scam, but the account is used for layering stolen funds.}
	\end{frame}
	
	\begin{frame}{A.12.1: Money Mule Account - Exam Scenario}
		\scenario{College Student Account with Multiple P2P Transfers}
		An account shows multiple incoming peer-to-peer transfers from dozens of different individuals under \$500 each, followed by immediate aggregated wire transfers to an offshore shell company account. The account holder is a college student with no verifiable income.
		
		\vspace{0.1cm}
		\textbf{What the Exam Tests:} Recognizing money mule account typology.
		
		\vspace{0.1cm}
		\textbf{How to Analyze:}
		\begin{enumerate}
			\item Multiple small incoming transfers from many senders = structured layering
			\item Immediate aggregation and wire to offshore shell company = mule network
			\item College student with no verifiable income = mule profile
			\item Pattern is classic money mule account for layering stolen funds
		\end{enumerate}
		
		\trap{Common Distractor: "The student is likely saving money from multiple jobs."}
		\correct{Correct Answer: The pattern is most likely a money mule account for layering stolen funds. The bank should file a SAR.}
	\end{frame}
	
	% ============================================================
	% SECTION 13: PREPAID CARDS AND CROSS-BORDER SMUGGLING
	% ============================================================
	
	\begin{frame}{A.13: Prepaid Cards as Money Laundering Vehicles}
		\begin{block}{Prepaid Card Vulnerabilities}
			Open-loop prepaid cards offer global liquidity, portability, and can bypass physical border declarations that apply to bulk cash.
		\end{block}
		\textbf{Prepaid Card Red Flags:}
		\begin{itemize}
			\item Purchasing multiple cards with cash at various retail locations
			\item Each card loaded just below reporting threshold (\$9,500 vs. \$10,000 CTR)
			\item Cards purchased over time to aggregate large value
			\item Traveler carrying numerous cards (50+ cards)
			\item Destination is high-risk or under-FinCEN-scrutiny jurisdiction
			\item Cards not linked to any bank account
		\end{itemize}
		\textbf{Structuring with Prepaid Cards:}
		Purchasing multiple prepaid cards with cash amounts just below the CTR threshold is a form of structuring - a federal crime independent of the underlying source of funds.
		
		\textbf{Border Reporting Evasion:}
		Prepaid cards are classified as monetary instruments under BSA regulations. Carrying multiple cards with aggregate value exceeding \$10,000 without declaration violates currency reporting requirements.
		
		\redflag{Prepaid cards can be used to structure cash and evade border currency reporting.}
	\end{frame}
	
	\begin{frame}{A.13.1: Prepaid Card Structuring - Exam Scenario}
		\scenario{Traveler Carrying 50 Prepaid Cards, Each Loaded with \$9,500}
		A traveler is stopped at customs carrying 50 prepaid debit cards, each loaded with \$9,500, purchased with cash at multiple retail locations over 30 days. The traveler is flying to a jurisdiction under FinCEN heightened scrutiny. Total value = \$475,000.
		
		\vspace{0.1cm}
		\textbf{What the Exam Tests:} Prepaid card structuring and border evasion.
		
		\vspace{0.1cm}
		\textbf{How to Analyze:}
		\begin{enumerate}
			\item Each card purchased at \$9,500 = just below \$10,000 CTR threshold
			\item Multiple cards purchased over time = structuring to avoid CTRs
			\item Cards are monetary instruments - aggregate value exceeds declaration threshold
			\item Traveler may have violated structuring laws AND currency reporting requirements
		\end{enumerate}
		
		\trap{Common Distractor: "Prepaid cards are retail products not subject to currency controls."}
		\correct{Correct Answer: The traveler may have violated structuring laws (purchasing at \$9,500) and evaded currency reporting requirements (prepaid cards are monetary instruments under BSA).}
	\end{frame}
	
	% ============================================================
	% SECTION 14: REAL ESTATE MONEY LAUNDERING
	% ============================================================
	
	\begin{frame}{A.14: Real Estate Money Laundering}
		\begin{block}{Real Estate Vulnerabilities}
			Real estate is an ideal integration vehicle - it holds value, generates legitimate income, and can be sold to convert illicit funds to clean proceeds.
		\end{block}
		\textbf{Real Estate Red Flags:}
		\begin{itemize}
			\item All-cash purchases through anonymous LLCs
			\item Purchase funds from pooled attorney trust accounts (IOLTA)
			\item Rapid LLC formation just before closing
			\item Attorney refusal to identify beneficial owner (citing privilege)
			\item Purchase price significantly above or below market value
			\item Rapid appreciation and sale with no legitimate explanation
			\item Funds sourced from offshore shell companies or secrecy havens
		\end{itemize}
		\textbf{FinCEN Geographic Targeting Orders (GTOs):}
		\begin{itemize}
			\item Require title insurance companies to identify beneficial owners behind all-cash purchases above threshold
			\item Target specific metropolitan areas identified as high-risk
			\item Apply to residential real estate transactions
		\end{itemize}
		\redflag{Anonymous LLC + IOLTA funds + attorney refusal = multiple red flags requiring EDD.}
	\end{frame}
	
	\begin{frame}{A.14.1: Real Estate Layering - Exam Scenario}
		\scenario{Luxury Condo Purchase Through Anonymous LLC}
		A \$6.5 million luxury condominium is purchased by a Delaware LLC formed 10 days before closing. Funds arrive from a law firm's IOLTA account. The LLC's registered agent is a corporate service provider known for forming anonymous companies. The law firm refuses to identify the beneficial owner, citing privilege.
		
		\vspace{0.1cm}
		\textbf{What the Exam Tests:} Recognizing real estate money laundering red flags.
		
		\vspace{0.1cm}
		\textbf{How to Analyze:}
		\begin{enumerate}
			\item Anonymous LLC (Delaware) with rapid formation = red flag
			\item IOLTA pooled funds = bank may not look through to individual client
			\item Attorney refusal to identify beneficial owner = significant red flag
			\item Combination of factors requires enhanced due diligence
		\end{enumerate}
		
		\trap{Common Distractor: "Licensed attorneys are exempt from AML scrutiny due to privilege."}
		\correct{Correct Answer: The combination of anonymous LLC, IOLTA funds, attorney refusal, and rapid LLC formation constitutes multiple red flags requiring EDD and potential SAR filing.}
	\end{frame}
	
	\begin{frame}{A.14.2: Real Estate Integration - Exam Scenario}
		\scenario{Hotel Purchase and Sale After 2 Years of Legitimate Operation}
		Structured funds (\$9,500 deposits over 24 months) are aggregated and used to purchase a beachfront resort for \$14 million in cash. The resort operates for 2 years, generating legitimate revenue. The criminal organization then sells the resort for \$18 million and invests proceeds in a stock portfolio.
		
		\vspace{0.1cm}
		\textbf{What the Exam Tests:} Recognizing integration through real estate.
		
		\vspace{0.1cm}
		\textbf{How to Analyze:}
		\begin{enumerate}
			\item Original \$9,500 deposits = placement
			\item Resort purchase = layering (converting cash to property)
			\item Resort operation and sale = INTEGRATION (legitimate rental income and sale proceeds)
			\item Rapid appreciation (\$14M to \$18M in 2 years) combined with anonymous LLC purchase warrants SAR filing
		\end{enumerate}
		
		\trap{Common Distractor: "The sale is between arms-length parties with a legitimate resort operation, so no SAR is needed."}
		\correct{Correct Answer: The sale of the resort and deposit of proceeds represents integration. The rapid appreciation combined with original anonymous LLC purchase and structured funding warrants a SAR filing.}
	\end{frame}
	
	% ============================================================
	% SECTION 15: LAWYER TRUST ACCOUNTS (IOLTA)
	% ============================================================
	
	\begin{frame}{A.15: IOLTA Account Vulnerabilities}
		\begin{block}{Why IOLTA Accounts Are Vulnerable}
			Banks often rely on the attorney's professional status without looking through to individual client sub-transactions. Multiple client funds are pooled together, obscuring the origin of specific funds.
		\end{block}
		\textbf{IOLTA Red Flags:}
		\begin{itemize}
			\item Large wire into IOLTA from offshore shell company
			\item Retainer amount disproportionate to the legal engagement
			\item Attorney has never met the client in person
			\item Rapid disbursement of funds after receipt
			\item Disbursement to offshore shell companies or other high-risk jurisdictions
			\item Attorney cannot explain the nature of the legal matter
		\end{itemize}
		\textbf{Exam Scenario Example:}
		A law firm receives \$5.5 million into its IOLTA account from a Cyprus shell company for a \$12,000 retainer. The partner has never met the client. The firm disburses \$5.2 million to a UAE shell company within 5 days, retaining \$300,000 as a "success fee."
		
		\trap{Common Distractor: "IOLTA accounts are exempt from AML reporting due to attorney-client privilege."}
		\correct{Correct Answer: The law firm's IOLTA functions as a blind pass-through vehicle. The disproportionate amount relative to retainer, offshore source, no face-to-face contact, and rapid disbursement are all red flags.}
	\end{frame}
	
	% ============================================================
	% SECTION 16: FREE TRADE ZONES (FTZs)
	% ============================================================
	
	\begin{frame}{A.16: Free Trade Zone Money Laundering}
		\begin{block}{FTZ Vulnerabilities}
			Relaxed regulatory oversight, minimal customs documentation inspections, and intense volume of re-exported goods create TBML opportunities.
		\end{block}
		\textbf{FTZ Red Flags:}
		\begin{itemize}
			\item Uniform invoice values regardless of cargo type
			\item Physical inspection reveals cargo different from declared goods
			\item Empty containers declared as containing goods (phantom shipping)
			\item Altered bills of lading after FTZ processing
			\item Nominee directors with no legitimate business role
			\item No verifiable manufacturing or logistics capability
		\end{itemize}
		\textbf{Exam Scenario Example:}
		A shipping company operating exclusively in an FTZ processed 340 containers declared as "industrial parts" with uniform invoice value of \$30,000. Inspection revealed one empty, one containing furniture, and one containing luxury electronics worth \$400,000. Bills of lading altered to show generic descriptions. Director is a legal nominee.
		
		\trap{Common Distractor: "FTZs are supervised by customs authorities, so no additional AML scrutiny is needed."}
		\correct{Correct Answer: The uniform invoice value regardless of cargo, empty containers (phantom shipping), mislabeling, and nominee director all constitute money laundering red flags requiring escalation.}
	\end{frame}
	
	% ============================================================
	% SECTION 17: CONCLUSION AND EXAM REMINDERS
	% ============================================================
	
	\begin{frame}{Domain A Summary: Top 30 Most Tested Concepts}
		\begin{multicols}{2}
			\begin{enumerate}
				\item Placement = introducing funds into system (structuring, cash deposits)
				\item Layering = obscuring the audit trail (wire transfers, shell companies)
				\item Integration = making funds appear legitimate (real estate, businesses)
				\item Over-invoicing transfers value TO exporter
				\item Under-invoicing transfers value TO importer
				\item Phantom shipping = no goods, only paper documentation
				\item Bearer shares = anonymous ownership (high risk)
				\item BMPE uses narco-dollars to pay U.S. commercial invoices
				\item Foreign PEPs require full EDD and senior management approval
				\item PEP family members also require EDD
				\item Human trafficking = ongoing exploitation (multiple small payments)
				\item Human smuggling = one-time border crossing (single large payment)
				\item NPO rapid cash withdrawal (80 percent+ in 48 hours) = TF red flag
				\item Privacy coins + mixers + cross-chain bridges = layering
				\item NFT wash trading = circular trading to inflate value
				\item Nested correspondent transactions = loss of visibility
				\item PTA requires correspondent to perform CDD on underlying accountholders
				\item Casino chip washing = placement + integration
				\item Hawala operators must register as MSBs (no cultural exemption)
				\item Money mules receive small transfers and wire to offshore shell companies
				\item Prepaid cards can structure cash and evade border reporting
				\item Real estate all-cash purchases through anonymous LLCs = red flag
				\item IOLTA accounts are blind pass-through vehicles
				\item FTZs enable phantom shipping and cargo mislabeling
				\item Structuring is a federal crime regardless of source of funds
				\item Self-laundering IS criminalized under FATF standards
				\item Trust UBOs include protector with veto power
				\item Dual-use goods to sanctioned countries = proliferation financing
				\item Terrorist financing has no minimum threshold
			\end{enumerate}
		\end{multicols}
	\end{frame}
	
	\begin{frame}{Critical Red Flags - Domain A Quick Reference}
		\begin{multicols}{2}
			\begin{itemize}
				\item \$9,500 cash deposits → Structuring
				\item 80 percent cash withdrawal in 48 hours for NPO → TF red flag
				\item Invoice price 80 percent below market → Under-invoicing (TBML)
				\item Multiple small payments to controller account → Human trafficking
				\item Anonymous LLC + IOLTA funds → Real estate layering
				\item Monero + mixer + Bitcoin → Privacy coin layering
				\item Undisclosed nested correspondent → Loss of visibility
				\item Chip purchase, minimal play, cash out → Casino chip wash
				\item Hawala with no MSB registration → Regulatory violation
				\item College student with mule pattern → Money mule account
				\item 50 prepaid cards at \$9,500 each → Structuring
				\item Bearer shares with no UBO disclosure → CDD failure
				\item Foreign PEP with unexplained wealth → EDD required
				\item Protector with veto power → Trust UBO
				\item Dual-use goods to sanctioned country → Proliferation financing
				\item \$10,000 TF threshold → Violates FATF Rec. 5
			\end{itemize}
		\end{multicols}
	\end{frame}
	
	\begin{frame}{Exam Day Reminders for Domain A}
		\begin{itemize}
			\item Domain A is 20 percent of exam
			\item Know the THREE stages of ML and be able to identify them in scenarios
			\item TBML techniques: over-invoicing (value to exporter) vs. under-invoicing (value to importer)
			\item Bearer shares are high-risk - immobilization does NOT eliminate CDD obligation
			\item Foreign PEPs require full EDD; family members also require EDD
			\item Human trafficking = ongoing exploitation (multiple small payments)
			\item Human smuggling = one-time border crossing (single large payment)
			\item NPO rapid cash withdrawal = terrorist financing red flag
			\item Privacy coins + mixers = high-risk layering pattern
			\item Nested correspondent accounts without disclosure = loss of visibility
			\item PTA requires correspondent to perform CDD on underlying accountholders
			\item Casino chip washing requires both CTR and SAR
			\item Hawala operators must register as MSBs - no cultural exemption
			\item Prepaid cards can be used for structuring and border evasion
			\item Real estate all-cash through anonymous LLCs = multiple red flags
			\item IOLTA accounts are blind pass-through vehicles
			\item Structuring is a federal crime regardless of source of funds
			\item Terrorist financing has NO minimum threshold
			\item Self-laundering IS criminalized under FATF Rec. 3
		\end{itemize}
		\vspace{0.3cm}
		\centering
		\greennote{Good luck on your CAMS exam!}
	\end{frame}
	
\end{document}